\section{Заключение}

В ходе выполнения дипломной работы были решены все поставленные задачи.

\begin{enumerate}
    \item Проведён анализ существующих методов автоматической векторизации, в частности подходов цикловой и SLP-векторизации. Выявлены преимущества и ограничения каждого метода, обоснован выбор SLP-подхода как наиболее подходящего для интеграции в компилятор Go.
    
    \item Исследована архитектура оптимизационного конвейера компилятора Go, включая SSA-представление, систему оптимизирующих проходов и особенности моделирования памяти через тип \texttt{TypeMem}. Установлены ключевые ограничения: отсутствие сложных анализов (псевдонимов, эволюции индуктивных переменных), требование быстрой компиляции и специфика рантайма (точный сборщик мусора, планировщик горутин).
    
    \item Разработан алгоритм выявления векторизуемых групп операций, включающий этапы поиска начальных пар смежных обращений к памяти, расширения групп по потоку данных (вверх и вниз), комбинирования, разбиения на компоненты связности, валидации (внешние использования, замкнутость операндов, когерентность линий) и планирования.
    
    \item Реализован прототип SLP-векторизатора в составе компилятора Go (пакет \texttt{cmd/compile/internal/ssa}, файл \texttt{slp.go}). Выполнена генерация 128-битных векторных инструкций для архитектуры ARM64 (NEON), реализована замена скалярных операций векторными с корректным обновлением цепочек памяти и вставкой извлечений для внешних использований.
    
    \item Проведена экспериментальная оценка разработанного решения на синтетических тестах, трассировщике лучей, промышленном приложении \texttt{etcd} и наборе \texttt{Sweet}. Подтверждена корректность преобразований и достигнуто ускорение до XX\% на вычислительно интенсивных фрагментах при незначительном увеличении времени компиляции (менее 1\%). Выявлены сценарии, где векторизация не применяется (различные состояния памяти, указатели в структурах, неполная ширина группы), что соответствует принятым компромиссам.
\end{enumerate}

Практическая значимость работы состоит в создании действующего прототипа SLP-векторизатора для компилятора Go, который:
\begin{itemize}
    \item демонстрирует принципиальную возможность интеграции SIMD-оптимизаций в стандартный компилятор;
    \item может служить основой для дальнейшего развития оптимизирующих возможностей \texttt{cmd/compile};
    \item предоставляет базу для расширения на другие архитектуры (x86-64, SVE) и поддержки более сложных векторизуемых шаблонов.
\end{itemize}

Дальнейшие направления развития включают:
\begin{itemize}
    \item введение упрощённой модели стоимости для более агрессивной векторизации;
    \item улучшение анализа памяти для распознавания независимых загрузок с разными \texttt{Memory}-аргументами;
    \item реализацию \texttt{shuffle}-инструкций для поддержки некогерентных линий;
    \item векторизацию редукций и частичных циклов.
\end{itemize}

Таким образом, цель работы — разработка и реализация прототипа механизма автоматической SLP-векторизации ациклических участков кода в компиляторе Go — достигнута. Полученные результаты подтверждают практическую пользу предложенного подхода и открывают путь к внедрению SIMD-оптимизаций в основной компилятор языка Go.